\documentclass{amsart}
\usepackage{amssymb}
\usepackage[hidelinks]{hyperref}

\newtheorem{theorem}{Theorem}[section]
\newtheorem{lemma}[theorem]{Lemma}
\newtheorem{proposition}[theorem]{Proposition}
\newtheorem{corollary}[theorem]{Corollary}
\theoremstyle{remark}
\newtheorem{remark}[theorem]{Remark}

\numberwithin{equation}{section}

\newcommand{\C}{\mathbb{C}}
\newcommand{\R}{\mathbb{R}}

\begin{document}

\title{Zeros of the independence polynomials of uniform linear hypercombs}

\author{Knowledge Sense}
\address{Knowledge Sense}
\email{math@knowledgesense.jp}

\date{9 October 2026}

\subjclass[2020]{Primary 05C31; Secondary 05A20, 05C65, 05C69, 30C15}

\keywords{Linear hypertree, hypercomb, strong independent set, independence polynomial, zeros of polynomials, log-concavity, unimodality}

\begin{abstract}
For $n\ge 1$ and $\ell\ge 3$ let $C_{n,\ell}$ be the $\ell$-uniform linear
hypercomb whose spine has $n$ edges. We prove that every zero $z$ of its
strong independence polynomial satisfies $|\arg z|>2\pi/3$. By a classical
criterion of Brenti, Royle and Wagner, the coefficients of this polynomial
are therefore strictly log-concave, and in particular unimodal. This
answers, for all $n$ and $\ell$, a question of Galvin and Sharpe, who proved
unimodality outside a window of length $O(\sqrt{n})$. Up to a linear factor,
the polynomial is the determinant of a tridiagonal matrix with two corner
corrections. Testing a unit kernel vector against this matrix reduces the
problem to a three-parameter family of quartic polynomials, all of whose
zeros we locate.
\end{abstract}

\maketitle

\section{Introduction}

Let $\mathcal{H}$ be a hypergraph. A set of vertices of $\mathcal{H}$ is a
\emph{strong independent set} if it meets every edge of $\mathcal{H}$ in at
most one vertex. If $i_k(\mathcal{H})$ denotes the number of strong
independent sets of size $k$, we write
\[
I(\mathcal{H};x)=\sum_{k\ge 0} i_k(\mathcal{H})\,x^k
\]
for the (strong) independence polynomial. For graphs this is the usual
independence polynomial, and whether every tree has a unimodal independence
sequence is a well-known open question of Alavi, Malde, Schwenk and
Erd\H{o}s~\cite{AMSE}. Galvin and Sharpe~\cite{GS} began the study of the
corresponding question for hypertrees. They restricted attention to
\emph{linear} hypertrees, in which two edges share at most one vertex, and
showed that linear hyperpaths have real-rooted, and linear hyperstars
log-concave, strong independence polynomials~\cite[Theorem~1.3]{GS}.

The next family considered in~\cite{GS} consists of hypercombs. Let
$n\ge 0$ and $\ell\ge 2$. The \emph{$\ell$-uniform linear hypercomb}
$C_{n,\ell}$ has spine vertices $u_0,\dots,u_n$ and $2n+1$ edges of size
$\ell$: for $1\le i\le n$ a \emph{spine edge} consisting of $u_{i-1}$, $u_i$
and $\ell-2$ further vertices, and for $0\le i\le n$ a \emph{tooth}
consisting of $u_i$ and $\ell-1$ further vertices. All vertices other than
$u_0,\dots,u_n$ lie in exactly one edge. For $n\ge 1$ this is the hypercomb
of~\cite[Section~1.1]{GS}; for $n=1$ the two teeth are attached at distinct
vertices of the unique spine edge, as required there. The hypergraph
$C_{0,\ell}$ is a single edge, and $I(C_{0,\ell};x)=1+\ell x$ is the initial
condition used in~\cite[Theorem~1.7]{GS}. For $\ell=2$ the hypercomb is the
centipede, whose independence polynomial is real-rooted~\cite{Zhu}.

For $\ell\ge 3$ the situation is different. Galvin and Sharpe found a linear
recurrence for $I(C_{n,\ell};x)$~\cite[Theorem~1.7]{GS} and observed by
computer that its coefficient sequence is unimodal for $n\le 500$ and
$\ell\le 12$. They proved that for every $\ell\ge 3$ there are constants
$K_\ell,L_\ell>0$ such that the sequence is weakly increasing up to the
index $K_\ell n-L_\ell\sqrt{n}$ and weakly decreasing from the index
$K_\ell n+L_\ell\sqrt{n}$~\cite[Corollary~1.8]{GS}; an earlier claim of
log-concavity for all $n$ and $\ell$ had been withdrawn~\cite[Note~1.9]{GS}.
They asked whether the window of length $O(\sqrt{n})$ can be shortened or
eliminated~\cite[Question~1.10(4)]{GS}. We eliminate it by locating the
zeros. Throughout, $\arg$ denotes the principal argument, with values in
$(-\pi,\pi]$.

\begin{theorem}\label{thm:main}
Let $n\ge 1$ and $\ell\ge 3$. Every zero $z$ of $I(C_{n,\ell};x)$ satisfies
$|\arg z|>2\pi/3$.
\end{theorem}

Since $I(C_{n,\ell};0)=1$, the theorem places all zeros in the open sector
of opening $2\pi/3$ around the negative real axis. A classical sufficient
condition of Brenti, Royle and Wagner~\cite[Proposition~1.2]{BRW}, in the
form recorded by Brown and Cameron~\cite[Corollary~2.3]{BC} and restated as
Proposition~\ref{prop:sector} below, converts this into information on the
coefficients.

\begin{corollary}\label{cor:lc}
Let $n\ge 1$ and $\ell\ge 3$, and write
$I(C_{n,\ell};x)=\sum_{k=0}^{2n+1}c_kx^k$. Then $c_k>0$ for
$0\le k\le 2n+1$, and
\[
c_k^2>c_{k-1}c_{k+1}\qquad(1\le k\le 2n).
\]
In particular the sequence $c_0,c_1,\dots,c_{2n+1}$ is log-concave and
unimodal.
\end{corollary}

Thus the window in~\cite[Corollary~1.8]{GS} is empty for every $n\ge 1$ and
$\ell\ge 3$, which settles~\cite[Question~1.10(4)]{GS}. The parts of
\cite[Question~1.10]{GS} concerning all linear hypertrees, or all uniform
linear hypertrees, remain open. We emphasize that the new statement is
Theorem~\ref{thm:main}; the passage from the location of zeros to
Corollary~\ref{cor:lc} uses only the classical criterion.

A recent preprint~\cite{Ast} treats a weighted family of hypercombs, defined
by the analogous recurrence with spine and tooth weights $s,t>0$; the uniform
hypercomb $C_{n,\ell}$ corresponds to $s=\ell-2$ and $t=\ell-1$. It is shown
there that for each fixed pair $s,t>0$ there is $N(s,t)$ such that the
coefficients are strictly log-concave for every $n\ge N(s,t)$, and the
question for all lengths in the unweighted case is left open. Corollary
\ref{cor:lc} applies only to the uniform weights, but covers every length
$n\ge 1$, and Theorem~\ref{thm:main} gives in addition the location of the
zeros. To our knowledge, Theorem~\ref{thm:main} has not appeared in the
literature.

\begin{remark}\label{rem:notreal}
The sector cannot be replaced by the negative real axis. For $n=5$ and
$\ell=3$ the coefficients $c_0,\dots,c_{11}$ of $I(C_{5,3};x)$ are
\[
1,\ 23,\ 220,\ 1156,\ 3716,\ 7700,\ 10564,\ 9716,\ 5984,\ 2400,\ 576,\ 64,
\]
as one checks from the recurrence~\eqref{eq:rec} below. Here
$c_{10}^2=331776<337920=\tfrac{11}{10}\cdot 2\cdot c_9c_{11}$, so Newton's
inequality $c_k^2\ge(1+1/k)(1+1/(d-k))c_{k-1}c_{k+1}$, valid for every
polynomial of degree $d$ with positive coefficients and only real zeros,
fails for $d=11$ and $k=10$. Hence $I(C_{5,3};x)$ has non-real zeros, in contrast with the
hyperpaths of~\cite[Theorem~1.3]{GS}.
\end{remark}

\subsection*{Method}
Up to the factor $1+(\ell-1)x$, the polynomial $I(C_{n,\ell};x)$ is the
determinant of an $n\times n$ tridiagonal matrix whose corner entries are
perturbed by $x$ (Lemmas~\ref{lem:cont} and~\ref{lem:det}). At a zero, a
unit kernel vector of this matrix produces two real numbers $\sigma$ and
$t$, the values of the quadratic forms of the corner and path matrices, with
$0\le\sigma\le 1$ and $|t|\le 2$, and $z$ is a zero of an explicit quartic
depending on $\sigma$ and $|t|$ (Lemma~\ref{lem:reduction}). The heart of
the proof is Lemma~\ref{lem:quartic}, which locates the zeros of all these
quartics. After the substitutions $w=-1/z$ and $w=a+v^2$, with $a=\ell-1$,
it becomes a statement about a quartic in $v$ with exactly two real zeros,
and Vieta's formulas determine the position of the remaining pair. No
realizability of the parameters is needed: the lemma covers the whole range
$0\le\sigma\le 1$ and $|t|\le 2$. The argument works verbatim with a real
parameter $a\ge 2$ in place of $\ell-1$ (Theorem~\ref{thm:general}).

\subsection*{Dependencies}
The proof of Theorem~\ref{thm:main} is self-contained. It uses the
recurrence of~\cite[Theorem~1.7]{GS}, which we rederive in
Lemma~\ref{lem:rec} to fix conventions. Corollary~\ref{cor:lc} uses in
addition the criterion of~\cite{BRW}.

\subsection*{Organization}
Section~\ref{sec:det} derives the recurrence and the determinant formula and
reduces Theorem~\ref{thm:main} to a family of quartics. Section~\ref{sec:quartic}
locates the zeros of these quartics. Section~\ref{sec:proofs} completes the
proofs.

\section{A determinant formula}\label{sec:det}

Let $a\ge 1$ be real, and put
\begin{equation}\label{eq:fg}
f=1+ax,\qquad g=1+(a-1)x,\qquad A=fg,\qquad B=xf.
\end{equation}
Define polynomials $C_n=C_n(x)$ by
\begin{equation}\label{eq:rec}
C_0=f+x,\qquad C_1=f(A+2x),\qquad C_n=AC_{n-1}+BC_{n-2}\quad(n\ge 2).
\end{equation}
Since $A(0)=1$ and $B(0)=0$, we have $C_n(0)=1$ for all $n\ge 0$.

\begin{lemma}\label{lem:rec}
Let $\ell\ge 2$ and $a=\ell-1$. Then $I(C_{n,\ell};x)=C_n(x)$ for every
$n\ge 0$.
\end{lemma}

\begin{proof}
Let $U_n$ and $V_n$ be the generating polynomials, by size, of the strong
independent sets of $C_{n,\ell}$ that do not contain, respectively contain,
the last spine vertex $u_n$. The single tooth of $C_{0,\ell}$ gives
$U_0=1+(\ell-1)x=f$ and $V_0=x$.

The hypercomb $C_{n+1,\ell}$ arises from $C_{n,\ell}$ by adding a spine edge
$\{u_n,u_{n+1}\}\cup X$ with $|X|=\ell-2$ and a tooth $\{u_{n+1}\}\cup Y$
with $|Y|=\ell-1$, where $X$, $Y$ and $\{u_{n+1}\}$ consist of new vertices.
A set $S$ of vertices of $C_{n+1,\ell}$ is a strong independent set if and
only if its restriction $S'$ to $C_{n,\ell}$ is one and the following holds.
If $u_{n+1}\notin S$, then $S$ contains at most one vertex of $Y$, and at
most one vertex of $X$, none if $u_n\in S'$. If $u_{n+1}\in S$, then
$u_n\notin S'$ and $S$ contains no vertex of $X\cup Y$. Since
$1+(\ell-2)x=g$, this gives
\[
U_{n+1}=f\,(gU_n+V_n),\qquad V_{n+1}=xU_n .
\]
Thus
\[
\begin{pmatrix} U_n\\ V_n\end{pmatrix}
=M^n\begin{pmatrix} U_0\\ V_0\end{pmatrix},\qquad
M=\begin{pmatrix} A & f\\ x & 0\end{pmatrix}.
\]
The matrix $M$ has trace $A$ and determinant $-B$, so $M^2=AM+BI$ by the
Cayley--Hamilton theorem, and $U_n+V_n$ satisfies the recurrence
in~\eqref{eq:rec} for $n\ge 2$. Finally $U_0+V_0=f+x$ and
$U_1+V_1=f(fg+x)+xf=f(A+2x)$.
\end{proof}

Since $A=fg$ and $B=xf$, the recurrence~\eqref{eq:rec} is
$C_n=f\,\bigl(gC_{n-1}+xC_{n-2}\bigr)$, which is the recurrence
of~\cite[Theorem~1.7]{GS}; one also checks that
$C_1=1+(3a+1)x+3a^2x^2+a^2(a-1)x^3$ agrees with the initial condition there.

Define $F_{-1}=0$, $F_0=1$ and $F_j=AF_{j-1}+BF_{j-2}$ for $j\ge 1$, so that
$F_1=A$.

\begin{lemma}\label{lem:cont}
For every $n\ge 1$,
\begin{equation}\label{eq:cont}
C_n=f\,\bigl(F_n+2xF_{n-1}+x^2F_{n-2}\bigr).
\end{equation}
\end{lemma}

\begin{proof}
Denote the bracket by $R_n$. Then $R_1=A+2x$, so~\eqref{eq:cont} holds for
$n=1$. For $n=2$,
\[
C_2=AC_1+BC_0=f(A^2+2xA)+xf(f+x)=f(A^2+B+2xA+x^2)=fR_2 .
\]
For $n\ge 3$ the indices $n$, $n-1$ and $n-2$ are positive, so each of
$F_n$, $F_{n-1}$, $F_{n-2}$ satisfies the defining recurrence and
$R_n=AR_{n-1}+BR_{n-2}$. Since $fR_n$ and $C_n$ satisfy the same recurrence
and agree for $n=1,2$, they agree for all $n\ge 1$.
\end{proof}

For $n\ge 2$ let $J_n$ be the adjacency matrix of the path with $n$ vertices,
that is, the $n\times n$ matrix with ones on the first super- and
subdiagonal and zeros elsewhere, and let $E_n$ be the diagonal matrix with
diagonal $(1,0,\dots,0,1)$.

\begin{lemma}\label{lem:det}
Let $n\ge 2$, $z\in\C$, and choose $\beta\in\C$ with $\beta^2=B(z)$. Then
the matrix
\[
M_n(z)=A(z)\,I_n+z\,E_n+i\beta\,J_n
\]
satisfies
$\det M_n(z)=F_n(z)+2zF_{n-1}(z)+z^2F_{n-2}(z)$.
\end{lemma}

\begin{proof}
Let $T_m$ be the $m\times m$ tridiagonal matrix with diagonal entries $A(z)$
and off-diagonal entries $i\beta$. Expansion along the last row gives
$\det T_m=A(z)\det T_{m-1}-(i\beta)^2\det T_{m-2}
=A(z)\det T_{m-1}+B(z)\det T_{m-2}$, with $\det T_0=1$ and
$\det T_1=A(z)$; hence $\det T_m=F_m(z)$. The matrix $M_n(z)$ arises from
$T_n$ by adding $z$ to the entries $(1,1)$ and $(n,n)$, which lie in
distinct rows because $n\ge 2$. By multilinearity of the determinant in
these two rows,
\[
\det M_n(z)=\det T_n+z\det T_{n-1}+z\det T_{n-1}+z^2\det T_{n-2},
\]
where the three minors are obtained by deleting row and column $1$, row and
column $n$, and both; for $n=2$ the last one is the empty matrix, with
determinant $1$. This is the asserted formula.
\end{proof}

\begin{lemma}\label{lem:reduction}
Let $n\ge 2$ and let $z\in\C$ satisfy $C_n(z)=0$. Then $z\ne 0$, and either
$z=-1/a$ or there are real numbers $\sigma\in[0,1]$ and $\tau\in[0,2]$ such
that
\begin{equation}\label{eq:quartic-z}
\bigl(A(z)+\sigma z\bigr)^2+\tau^2 z f(z)=0 .
\end{equation}
\end{lemma}

\begin{proof}
We have $z\ne 0$ because $C_n(0)=1$. Suppose $z\ne -1/a$, so that
$f(z)\ne 0$. By Lemmas~\ref{lem:cont} and~\ref{lem:det}, $\det M_n(z)=0$.
Let $v\in\C^n$ satisfy $M_n(z)v=0$ and $v^*v=1$, and put
\[
\sigma=v^*E_nv=|v_1|^2+|v_n|^2,\qquad
t=v^*J_nv=2\operatorname{Re}\sum_{j=1}^{n-1}\overline{v_j}\,v_{j+1}.
\]
Then $\sigma\in[0,1]$, $t$ is real, and
$|t|\le\sum_{j=1}^{n-1}\bigl(|v_j|^2+|v_{j+1}|^2\bigr)\le 2$. Moreover
\[
0=v^*M_n(z)v=A(z)+\sigma z+i\beta t .
\]
Hence $(A(z)+\sigma z)^2=-\beta^2t^2=-t^2zf(z)$, and~\eqref{eq:quartic-z}
holds with $\tau=|t|$.
\end{proof}

\section{A family of quartics}\label{sec:quartic}

\begin{lemma}\label{lem:quartic}
Let $a\ge 2$, $\sigma\in[0,1]$ and $\tau\in[0,2]$ be real. Every zero $z$
of
\[
\Phi(z)=\bigl((1+az)(1+(a-1)z)+\sigma z\bigr)^2+\tau^2z(1+az)
\]
satisfies $|\arg z|>2\pi/3$.
\end{lemma}

\begin{proof}
If $\tau=0$, then $\Phi$ is the square of
$1+(2a-1+\sigma)z+a(a-1)z^2$. This quadratic has positive coefficients and
discriminant $(2a-1+\sigma)^2-4a(a-1)=1+2\sigma(2a-1)+\sigma^2>0$, so its
zeros are negative real numbers. We assume from now on that $\tau>0$.

\emph{Step 1: two substitutions.} Since $\Phi(0)=1$, every zero $z$ of
$\Phi$ is nonzero, and we put $w=-1/z$. From $w(1+az)=w-a$,
$w(1+(a-1)z)=w-a+1$ and $wz=-1$ we obtain
\[
w^4\,\Phi(-1/w)=D(w)^2-\tau^2w^2(w-a),\qquad
D(w)=(w-a)(w-a+1)-\sigma w .
\]
If $|\arg w|<\pi/3$, then $|\arg z|=\pi-|\arg w|>2\pi/3$. It therefore
suffices to show that every zero $w$ of $D(w)^2-\tau^2w^2(w-a)$ satisfies
$|\arg w|<\pi/3$. Let $w$ be such a zero, and choose $v\in\C$ with
$v^2=w-a$. Then $D(w)^2=(\tau wv)^2$, and after replacing $v$ by $-v$ if
necessary, $D(w)=\tau wv$. Substituting $w=a+v^2$, so that
$D(w)=v^2(v^2+1)-\sigma(a+v^2)$, we find that $v$ is a zero of
\[
\Psi(v)=v^4-\tau v^3+(1-\sigma)v^2-a\tau v-a\sigma .
\]
Consequently it suffices to prove that
\begin{equation}\label{eq:goal}
|\arg(a+v^2)|<\pi/3\qquad\text{for every zero $v$ of $\Psi$.}
\end{equation}

\emph{Step 2: the real zeros of $\Psi$.} For $0<v\le\tau$ we have
$v^3(v-\tau)\le 0$ and $(1-\sigma)v^2\le(1-\sigma)\tau v$, hence
\[
\Psi(v)\le(1-\sigma-a)\tau v-a\sigma<0 ,
\]
because $a>1\ge 1-\sigma$. For $v\ge\tau$ we have
$\Psi''(v)=6v(2v-\tau)+2(1-\sigma)>0$. As $\Psi(\tau)<0$ and
$\Psi(v)\to\infty$, the strictly convex function $\Psi$ has exactly one zero
$r$ in $(\tau,\infty)$, and $\Psi'(r)>0$. So $r>\tau$ is the only positive
zero of $\Psi$, it is simple, and $\Psi<0$ on $(0,r)$. For $p\ge 0$,
\[
\Psi(-p)=p^4+\tau p^3+(1-\sigma)p^2+a\tau p-a\sigma
\]
is strictly increasing in $p$, with derivative at least $a\tau>0$, starts
at $-a\sigma\le 0$ and tends to infinity. Hence $\Psi$ has exactly one zero
$-p$ with $p\ge 0$, and it is simple. Since $\Psi$ has real coefficients,
its remaining zeros form a pair $\xi,\overline{\xi}$ with
$\operatorname{Im}\xi\ne 0$. Moreover $p<r$: this is clear if $p=0$, and if
$p>0$ then $\Psi(p)=\Psi(-p)-2\tau p(p^2+a)<0$, so $p\in(0,r)$.

For the real zeros, $a+r^2$ and $a+p^2$ are positive, so~\eqref{eq:goal}
holds. The zero $\overline{\xi}$ yields the complex conjugate of $a+\xi^2$,
so it remains to show that $|\arg(a+\xi^2)|<\pi/3$.

\emph{Step 3: Vieta's formulas.} Put
\[
\delta=r-p-\tau,\qquad Q=|\xi|^2>0,\qquad H=rp\ge 0 .
\]
Comparing the coefficients of $v^3$ in
$\Psi(v)=(v-r)(v+p)(v-\xi)(v-\overline{\xi})$ gives
$r-p+2\operatorname{Re}\xi=\tau$, that is, $2\operatorname{Re}\xi=-\delta$.
Hence $(v-\xi)(v-\overline{\xi})=v^2+\delta v+Q$, while
$(v-r)(v+p)=v^2-(\tau+\delta)v-H$. Comparing the remaining coefficients, we
obtain
\begin{align}
Q&=1-\sigma+H+\delta(\tau+\delta),\label{eq:v2}\\
a&=Q+\delta\,\frac{Q+H}{\tau},\label{eq:v1}\\
HQ&=a\sigma .\label{eq:v0}
\end{align}

\emph{Step 4: $\delta>0$.} Suppose that $\delta\le 0$. Since $r>p$, we have
$\delta>-\tau$, and therefore $\delta(\tau+\delta)\le 0$. By~\eqref{eq:v1},
$a\le Q$, so~\eqref{eq:v0} gives $H=a\sigma/Q\le\sigma$. Then~\eqref{eq:v2}
gives $Q\le 1-\sigma+\sigma=1$, which contradicts $Q\ge a\ge 2$.

\emph{Step 5: the sector.} Put $K=(Q+H)/\tau$. Then $K>0$, $K\ge Q/2$
because $H\ge 0$ and $\tau\le 2$, and $a=Q+\delta K$ by~\eqref{eq:v1}.
Write $\xi=-\delta/2+iy$ with $y\ne 0$, so that $Q=\delta^2/4+y^2$. Then
\[
a+\xi^2=Q+\delta K+\frac{\delta^2}{4}-y^2-i\delta y
=\delta\Bigl(K+\frac{\delta}{2}\Bigr)-i\delta y .
\]
The real part is positive by Step~4. Hence $|\arg(a+\xi^2)|<\pi/3$ if and
only if $\delta^2y^2<3\delta^2(K+\delta/2)^2$, that is,
\begin{equation}\label{eq:final}
Q<3K^2+3K\delta+\delta^2 .
\end{equation}
If $Q\ge 4/3$, then $3K^2\ge 3Q^2/4\ge Q$, and~\eqref{eq:final} holds
because $3K\delta+\delta^2>0$. If $Q<4/3$, then
$3K\delta=3(a-Q)>3(2-4/3)=2>Q$, and~\eqref{eq:final} holds again. This
proves~\eqref{eq:goal} and the lemma.
\end{proof}

\section{Proofs of the main results}\label{sec:proofs}

\begin{theorem}\label{thm:general}
Let $a\ge 2$ be real and let $C_n$ be defined by~\eqref{eq:fg}
and~\eqref{eq:rec}. For every $n\ge 0$, every zero $z$ of $C_n$ satisfies
$|\arg z|>2\pi/3$.
\end{theorem}

\begin{proof}
The only zero of $C_0=1+(a+1)x$ is $-1/(a+1)$. For $n=1$,
$C_1=f\,(A+2x)$, where $f$ has the zero $-1/a$ and
$A+2x=1+(2a+1)x+a(a-1)x^2$ has positive coefficients and discriminant
$8a+1>0$; so all zeros of $C_1$ are negative real numbers. Let $n\ge 2$ and
let $z$ be a zero of $C_n$. By Lemma~\ref{lem:reduction}, either $z=-1/a$,
or~\eqref{eq:quartic-z} holds for some $\sigma\in[0,1]$ and
$\tau\in[0,2]$. In the latter case $z$ is a zero of the polynomial $\Phi$
of Lemma~\ref{lem:quartic}, because $A(z)=(1+az)(1+(a-1)z)$ and
$zf(z)=z(1+az)$; hence $|\arg z|>2\pi/3$.
\end{proof}

\begin{proof}[Proof of Theorem~\ref{thm:main}]
Apply Lemma~\ref{lem:rec} and Theorem~\ref{thm:general} with
$a=\ell-1\ge 2$.
\end{proof}

The passage to coefficients rests on the following classical criterion.
Brenti, Royle and Wagner~\cite[Proposition~1.2]{BRW} state it for zeros in
the sector $|\arg z|<\pi/3$; replacing $p(x)$ by $p(-x)$ changes the sign of
every second coefficient and does not affect the inequalities, which gives
the form below~\cite[Corollary~2.3]{BC}.

\begin{proposition}[\cite{BRW,BC}]\label{prop:sector}
Let $p(x)=\sum_{k=0}^d b_kx^k\in\R[x]$ have degree $d$ and $p(0)\ne 0$,
and suppose that every zero $z$ of $p$ satisfies $|\arg z|>2\pi/3$. Then the coefficients
$b_0,\dots,b_d$ are all positive or all negative, and
$b_k^2>b_{k-1}b_{k+1}$ for $1\le k\le d-1$.
\end{proposition}

\begin{proof}[Proof of Corollary~\ref{cor:lc}]
The $2n+1$ edges of $C_{n,\ell}$ cover all its vertices, so a strong
independent set has at most $2n+1$ elements. Choosing one vertex outside the
spine in every edge, which is possible because $\ell\ge 3$, gives one of
size $2n+1$. Hence $I(C_{n,\ell};x)$ has degree $2n+1$. By
Theorem~\ref{thm:main} and Proposition~\ref{prop:sector}, its coefficients
have a common sign, which is positive because $c_0=1$, and
$c_k^2>c_{k-1}c_{k+1}$ for $1\le k\le 2n$. Therefore the ratios
$c_k/c_{k-1}$, $1\le k\le 2n+1$, are strictly decreasing, and the sequence
is unimodal.
\end{proof}

\begin{remark}\label{rem:weak}
For the non-strict inequalities $c_k^2\ge c_{k-1}c_{k+1}$ one can also argue
directly. By Theorem~\ref{thm:main}, $I(C_{n,\ell};x)$ is a product of
factors $1+x/\rho$ and $1-(2\cos\theta)x/\rho+x^2/\rho^2$ with $\rho>0$ and
$2\pi/3<|\theta|<\pi$. Since $-2\cos\theta\in(1,2)$, every factor has
positive and log-concave coefficients, and the convolution of positive
log-concave sequences without internal zeros is again of this kind; see,
for example,~\cite{Karlin}.
\end{remark}

\begin{remark}
By Theorem~\ref{thm:general} and Proposition~\ref{prop:sector}, the
conclusion of Corollary~\ref{cor:lc} holds for the polynomials $C_n$,
$n\ge 1$, for every real $a\ge 2$. The assumption $a\ge 2$ is used in Steps
4 and 5 of the proof of Lemma~\ref{lem:quartic}. The case $a=1$ corresponds
to the centipede, whose independence polynomial is real-rooted~\cite{Zhu}.
\end{remark}

\section*{Use of AI-assisted tools}
AI-assisted tools were used in the research and preparation of this
manuscript.

\begin{thebibliography}{9}

\bibitem{AMSE}
Y.~Alavi, P.~J. Malde, A.~J. Schwenk, and P.~Erd\H{o}s, \emph{The vertex
independence sequence of a graph is not constrained}, Congr. Numer.
\textbf{58} (1987), 15--23.

\bibitem{BRW}
F.~Brenti, G.~F. Royle, and D.~G. Wagner, \emph{Location of zeros of
chromatic and related polynomials of graphs}, Canad. J. Math. \textbf{46}
(1994), no.~1, 55--80.

\bibitem{BC}
J.~I. Brown and B.~Cameron, \emph{On the unimodality of independence
polynomials of very well-covered graphs}, Discrete Math. \textbf{341}
(2018), no.~4, 1138--1143; arXiv:1709.08236.

\bibitem{GS}
D.~Galvin and C.~Sharpe, \emph{Independent set sequence of some linear
hypertrees}, Australas. J. Combin. \textbf{94} (2026), no.~1, 93--121;
arXiv:2409.15555.

\bibitem{Karlin}
S.~Karlin, \emph{Total Positivity, Vol.~I}, Stanford University Press,
Stanford, 1968.

\bibitem{Ast}
OpenAI Codex GPT-6-Astra, \emph{Eventual strict log-concavity of weighted
linear hypercombs}, preprint (2026), aiXiv:260910.000002.

\bibitem{Zhu}
Z.-F. Zhu, \emph{The unimodality of independence polynomials of some
graphs}, Australas. J. Combin. \textbf{38} (2007), 27--33.

\end{thebibliography}

\end{document}
